\section{The proper quadratic quotient in the full ternary-rich cone}
Let $\lambda\in(1,2)$ satisfy $\lambda^3=\lambda^2+\lambda+1$,
and set $a=\lambda-1$, $b=\lambda^{-1}=\lambda^2-\lambda-1$,
$B(z)=1+az+bz^2$. Write $[t]=1+\cdots+z^{t-1}$.
For $E=[2]^u[3]^r$, let
\[
                         E=BJ+R_0+R_1z                 \tag{1}
\]
be its actual polynomial division by B. Unlike the preceding forward
inverse theorem, J is a polynomial quotient, not a truncated FPS.

\begin{theorem}[All rich counts, without a binary-count bound]
For EVERY $u\ge4$, $r\ge2u$, J is strictly positive and globally
log-concave. Its two remainder coefficients satisfy
$|R_0|,|R_1|\le1/4$ and the exact identity
\[
 bR_0^2-aR_0R_1+R_1^2
       =b\sigma^u b^r,
 \quad\sigma=1+2\lambda-\lambda^2=a^2/b>1.         \tag{2}
\]
\end{theorem}
In fact the same positive/global-LC conclusion holds for every
$r\ge18,u\ge0$ such that $b\sigma^u b^r\le1/64$.
This is an unbounded norm cone, not an experimentally tested count range.

\paragraph{Two exact norm recurrences, for both boundary pairs.}
Put $\mathcal R(x,y)=bx^2+axy+y^2$. Define
\[
 H=\frac{E}{1+(a/b)z+(1/b)z^2},\quad D=u+2r,
 \quad(e_0,e_1)=(H_D,H_{D-1}).
\]
Reversal of (1) gives
$J^{\rm rev}=b^{-1}[H]_{0..D-2}$. Its sublist beginning at
index two is therefore the reversal of $b^{-1}[H]_{0..D-4}$.
This is why four coefficients, rather than the preceding forward
kernel's two, are initially discarded in the present argument.

On binary multiplication the H terminal pair changes by
\[
 (x,y)\mapsto((1-\lambda a)x-\lambda y,x+y),
                        \quad\mathcal R_{\rm new}=\sigma\mathcal R.
\]
On ternary multiplication it changes by
\[
              (x,y)\mapsto(y,-bx-ay),
                        \quad\mathcal R_{\rm new}=b\mathcal R.
\]
The latter simplification uses the complete cubic identities
$1-\lambda a+(\lambda a)^2-\lambda=0$ and
$\lambda(\lambda a-1)=1$.
For the signed remainder pair $(R_0,-R_1)$ the corresponding maps are
\[
 (x,y)\mapsto(x+\lambda y,-x+(1-\lambda a)y),
 \qquad (x,y)\mapsto(-ax-y,bx),
\]
with the SAME factors sigma and b. Both base pairs for E=1 are (1,0)
and have norm b. Thus complete two-count induction proves
\[
 \mathcal R(e_0,e_1)=\mathcal R(R_0,-R_1)=b\sigma^u b^r. \tag{3}
\]
Both phase variables and both remainder coefficients are retained.
The two norm reports verify BOTH terminal/remainder maps; a third
report formalizes the entire two-natural-count norm induction.

\paragraph{A uniform small norm throughout the rich cone.}
Here $b\ge1/2$, $a^2\le3/4$, and
\[
 \mathcal R(x,y)-\tfrac14(x^2+y^2)
       =(b-\tfrac12)x^2+(\tfrac34-a^2)y^2+(x/2+ay)^2\ge0.
\]
Hence norm at most 1/64 gives the full quarter box for either pair.
The exact relations $a^2b=\lambda^2-3<2/5$, $0<b<3/5$ imply,
for $u\ge4,r=2u+k$,
\[
 b\sigma^u b^r=b(a^2b)^u b^k
              \le\tfrac35(\tfrac25)^4
                    =48/3125<1/64.                       \tag{4}
\]
Strong coercivity/quarter-box and the all-u/all-k geometric payment
are separate scoped Lean reports. The fixed algebraic root bounds
used here are among the complete scalar certificate.

\paragraph{Every uncovered reverse-prefix read is paid.}
Let $w=u+r\ge18$, and write $x_k=H_{D-k}$. Palindromicity gives
\[
 x_0=e_0,\ x_1=e_1,\qquad
 x_k=bE_{k-2}-ax_{k-1}-bx_{k-2}\quad(2\le k\le7),          \tag{5}
\]
where every needed coefficient is the explicit polynomial
\[
               E_t=\sum_{j\le t/2}{r\choose j}{w-j\choose t-2j}.
\]
For an old positive LC prefix through D-4, ordinary interval
convolution covers all new centres through D-5. Binary extension
leaves only one new value and one interior minor:
\[
 x_3+x_4,\qquad (x_4+x_5)^2-(x_3+x_4)(x_5+x_6).
\]
Ternary extension leaves two new values and two minors:
\[
\begin{gathered}
 x_3+x_4+x_5,\quad x_2+x_3+x_4,\\
 (x_4+x_5+x_6)^2-(x_3+x_4+x_5)(x_5+x_6+x_7),\\
 (x_3+x_4+x_5)^2-(x_2+x_3+x_4)(x_4+x_5+x_6).
\end{gathered}
\]
All SIX expressions are strictly positive for ALL real w>=18,
0<=r<=w and the entire terminal quarter box. Substitute
$w=18+X$, $r=\theta w$, $e_i=-1/4+\eta_i/2$ and expand each
coefficient of X in the tensor Bernstein basis on the unit cube.
The respective coefficient counts are 99,143,99 for the minors and
11,13,11 for the values: 376 in total. Every coefficient is
nonnegative, and all at $X^0$ are positive. The complete list and
exact sign certificate are supplied. Independence is checked by
complete tensor interpolation, using degree bound (8,4,2,2), not
by a bounded parameter test.

\paragraph{A single pure entrance proves the entire norm cone.}
First fix u=0. The one explicit r=18 polynomial has positive LC
J[2..], equivalently positive LC H[0..D-4]. Its norm b^19 is below
1/64. Further ternary multiplication contracts the norm, so (5)
and the six corner inequalities prove this property for every r>=18.

Now fix any r>=18 and any target u with $b\sigma^u b^r\le1/64$.
Starting from that pure polynomial, add u binary factors. Since
sigma>1, every previous norm is at most the TARGET norm; every
old terminal pair is consequently in the required box. Binary
convolution and its complete endpoint payments prove LC of J[2..]
at EVERY intermediate count. The formal increasing-norm predicate
induction includes all natural u. No unbounded diagonal is assumed
as a finite entrance.

\paragraph{Pay BOTH missing head minors and every low coefficient.}
The remainder box from (3) gives
\[
 \begin{aligned}
 J_0&=1-R_0,& J_1&=w-aJ_0-R_1,\\
 J_2&={w\choose2}+r-aJ_1-bJ_0,&
 J_3&={w\choose3}+r(w-1)-aJ_2-bJ_1.
 \end{aligned}
\]
All three values J0,J1,J2 and BOTH minors
$J_1^2-J_0J_2$, $J_2^2-J_1J_3$ are positive whenever w>=12,
0<=r<=w, |R0|,|R1|<=1/4. The complete analogous Bernstein
certificate has 2,5,11,28,56 coefficients respectively: 102 in total.
This explicitly pays centre TWO; LC of J[2..] alone covers only
centres from three and cannot replace this payment.
It proves global positivity and LC throughout the norm cone.

By (4), this proves every rich case with r>=18. The remaining rich
cases are EXACTLY thirty: u=4,r=8..17; u=5,r=10..17;
u=6,r=12..17; u=7,r=14..17; u=8,r=16..17. All their complete
quotients and minors, together with the one pure entrance, form
31 explicit contexts with 993 retained coefficients and 931 minors.
They are independently reconstructed in the different Q(a) basis.
Their signs, all 478 Bernstein coefficients and the fixed root
obligations use 2350 distinct field triples in 25 ordinary Lean
integer chunks, represented by ONE aggregate report.

\paragraph{Scope.}
Seven NEW reports comprise six full norm/coercivity/rich-bound/
increasing-cone kernels and the complete finite scalar certificate.
The earlier cubic-root interval and finite Bernstein-sum kernels
are reused, not recounted. Sequence/quotient identifications,
LC convolution, the independently verified polynomial identities
and the full norm-cone induction are the written assembly above,
not an end-to-end Lean formalization of composition generation.
Positive-primitive LC of J remains a SEPARATE unproved candidate;
it is not inferred from ordinary LC. The full actual D profile
contains additional source terms. This quotient theorem alone
does not prove the actual strong diagonal for u>=33, general C,
full A2 or ORIGINAL undilated coefficient unimodality.
